\documentclass[10pt,a4paper]{article}

\usepackage[utf8]{inputenc}
\usepackage[T1]{fontenc}
\usepackage{lmodern}
\usepackage[french]{babel}
\usepackage[margin=1.4cm]{geometry}
\usepackage{microtype}
\usepackage{multicol}
\usepackage{enumitem}
\usepackage[table]{xcolor}
\usepackage[most]{tcolorbox}
\usepackage{amsmath}

\definecolor{navy}{HTML}{1F3864}
\definecolor{lightblue}{HTML}{D6E0F0}
\definecolor{warm}{HTML}{C55A11}

\newtcolorbox{alerte}[1][]{colback=orange!8,colframe=orange!70!black,boxrule=0.6pt,
  arc=2pt,left=6pt,right=6pt,top=4pt,bottom=4pt,#1}
\newtcolorbox{reflexe}[1][]{colback=lightblue!50,colframe=navy,boxrule=0.8pt,
  arc=2pt,left=6pt,right=6pt,top=4pt,bottom=4pt,#1}

\newcommand{\slide}[2]{%
  \medskip\noindent{\color{navy}\bfseries #1.}\ \textcolor{navy}{\footnotesize #2}\par\nobreak\smallskip}
\newcommand{\dire}[1]{\noindent\hangindent=1em\hangafter=1 #1\par}
\newcommand{\n}[1]{\textbf{#1}}

\setlength{\parindent}{0pt}
\setlength{\columnsep}{16pt}
\pagestyle{empty}

\begin{document}

{\color{navy}\Large\bfseries Aide-mémoire, point tuteur du 31 juillet 2026}\\[2pt]
{\footnotesize Kélian Kaddouri. À garder sous les yeux. Les chiffres en gras sont ceux qu'il ne faut pas se tromper.}

\vspace{4pt}
\begin{reflexe}
\textbf{La phrase à placer deux fois, minimum} (slide 8, puis en clôture)~:\\
\textbf{\og Le niveau bouge dans la bande, l'ordre des piliers ne bouge pas. C'est ça, mon résultat.\fg{}}\\[3pt]
{\footnotesize Sans elle, la liste des incertitudes passe pour du flou. Avec elle, elle passe pour de la rigueur.}
\end{reflexe}

\vspace{2pt}
\begin{multicols}{2}
\footnotesize

\slide{1. Titre}{}
Point d'avancement. Les quatre pistes du 24/07, et ce qu'elles ont ouvert.

\slide{2. Depuis le dernier point}{annoncer la demande}
\dire{$\bullet$ Quatre pistes données, quatre traitées.}
\dire{$\bullet$ Deux ont ouvert des sujets plus larges $=$ 2\textsuperscript{e} moitié.}
\dire{$\bullet$ \n{\og À la fin, j'ai un service à te demander.\fg}}

\slide{3. Comptage}{N2}
\dire{$\bullet$ Trois mesures de surdispersion jamais confrontées.}
\dire{$\bullet$ Deux sont la même, coïncident à \n{5\,\%}.}
\dire{$\bullet$ À variance égale, la queue bouge de \n{11\,\%}.}
\dire{$\bullet$ Le comptage déplace \n{la VaR}, pas la moyenne.}
\dire{$\bullet$ Aveu volontaire~: série détrendée $=$ surdispersion \n{3 fois} plus faible. \og conservateur, pas calibré\fg.}

\slide{4. Seuil P4}{}
\dire{$\bullet$ Ta question en cachait deux.}
\dire{$\bullet$ \n{Marge} (P4 bascule seul)~: normale inoffensive, elle sert d'unité. La vraie queue est dans la sévérité.}
\dire{$\bullet$ \n{Dépendance} (combien ensemble)~: gaussien $=$ dépendance de queue \n{nulle}, contredit le Gumbel du mémoire.}
\dire{$\bullet$ Donc j'encadre au lieu de trancher.}
\dire{$\bullet$ Co-extrême~: \n{0,17 $\to$ 0,39}. C'est MOVEit, que le gaussien rate.}

\slide{5. $W=S+A$}{Z11 -- le meilleur moment}
\dire{$\bullet$ \n{\og Ton intuition était juste, et c'est une identité, pas une analogie.\fg}}
\dire{$\bullet$ On coupe en deux~: \n{qui est lié à qui}, et \n{dans quel sens}.}
\dire{$\bullet$ Renverse le temps~: la 1\textsuperscript{re} ne bouge pas, la 2\textsuperscript{e} s'inverse.}
\dire{$\bullet$ La fluctuation martingale ne dépend que de la 1\textsuperscript{re}.}
\dire{$\bullet$ Donc une donnée annuelle, sans ordre, ne peut pas voir la direction.}
\dire{$\bullet$ \n{\og Mon placebo à $-0{,}33$ n'est pas un échec, c'est le théorème.\fg}}

\slide{6. KPI}{Z13}
\dire{$\bullet$ Un score de criticité ne déclenche aucune décision.}
\dire{$\bullet$ Chaque indicateur branché sur un canal, capital chiffré derrière.}
\dire{$\bullet$ Le point n'est pas le montant, c'est \n{la concentration}.}
\dire{$\bullet$ Donc on priorise au lieu d'étaler. Défendable devant une DAF.}

\slide{7. Caroline}{}
\dire{$\bullet$ \og Un quantile d'un objet incertain est mauvais.\fg{} Elle a raison.}
\dire{$\bullet$ Mesurable~: le seul IC vaut un facteur \n{2,6}.}
\dire{$\bullet$ La VaR est imposée par le cadre, je ne peux pas l'abandonner.}
\dire{$\bullet$ Donc je rends l'incertitude visible~: \n{paramètre, identification, modèle}.}
\dire{$\bullet$ \n{\og Sa remarque ne casse pas l'approche, elle la justifie.\fg}}

\slide{8. Trois étages}{J4 -- placer la phrase}
\dire{$\bullet$ Dépendance~: facteur \n{1,8}. Famille de queue~: facteur \n{5,5}.}
\dire{$\bullet$ Ma cascade tombe entre gaussien et Gumbel.}
\dire{$\bullet$ \n{\og Le niveau bouge, l'ordre des piliers ne bouge pas.\fg}}

\slide{9. Validation}{J3}
\dire{$\bullet$ \n{\og Et tout n'est pas incertain.\fg}}
\dire{$\bullet$ Tests non rejetés, \n{$p=0{,}87$}, bootstrap paramétrique.}
\dire{$\bullet$ Négative binomiale bat le Poisson sans discussion.}
\dire{$\bullet$ Les \n{formes} sont validées, le \n{niveau} reste une bande.}

\slide{10. Direction}{}
\dire{$\bullet$ Aucune base ne mesure le sens. Mais chaque grand incident a son enquête publique.}
\dire{$\bullet$ \n{7 rapports}, \n{17 enchaînements}, protocole écrit \emph{avant} de coder.}
\dire{$\bullet$ Placebo cette fois franchi~: \n{$z=+3{,}9$} contre $-0{,}33$.}
\dire{$\bullet$ Aucun enchaînement à contre-sens.}
\dire{$\bullet$ Gouvernance jamais cible, incidents toujours aboutissement.}
\dire{$\bullet$ \n{\og La structure du classeur, retrouvée par un chemin indépendant.\fg}}

\slide{11. Valeur d'information}{Z18}
\dire{$\bullet$ Ces 7 rapports ont levé \n{66\,\%} de l'ambiguïté.}
\dire{$\bullet$ Et ça se \n{hiérarchise}~: gouvernance $\to$ incidents en lèverait \n{29\,\%} seule, une autre \n{rien}.}
\dire{$\bullet$ La recommandation au régulateur cesse d'être générique.}
\dire{$\bullet$ Et la plus précieuse est celle que les rapports ne couvrent pas.}

\slide{12. La demande}{le moment clé}
\dire{$\bullet$ Le tableau~: structure \n{revendiquée}, sens de chaque lien \n{non revendiqué} (2 paires sur 6), force \n{hors d'atteinte}.}
\dire{$\bullet$ \n{\og Ce dépouillement, je l'ai fait seul. Et une enquête cherche un responsable, donc elle conclut presque toujours à la gouvernance.\fg}}
\dire{$\bullet$ \n{\og Il me faut un second codeur qui ne connaisse pas mes conclusions. Donc pas toi, et c'est précisément le point.\fg}}
\dire{$\bullet$ 45 min, aucune maths.}
\dire{$\bullet$ \n{\og Qui verrais-tu chez Nexialog ?\fg} $\leftarrow$ demander un \n{nom}, pas un accord.}
\dire{$\bullet$ Puis~: relecture praticien, 30 min, 7 phrases à contredire. \n{Toi}, et \n{Mehdi} à son retour.}

\slide{13. P4 sous-process}{Z14 + Z16}
\dire{$\bullet$ Taxonomie fine du tiers inexistante publiquement $\Rightarrow$ le registre reste nécessaire.}
\dire{$\bullet$ Ancrés sur donnée européenne officielle~: origination \n{29\,\%}, concentration.}
\dire{$\bullet$ Si le facteur tiers est un processus, le seuil devient un \n{temps d'arrêt}.}
\dire{$\bullet$ \og Déclenchement\fg{} prend un sens littéral, et rejoint la 3\textsuperscript{e} piste.}

\slide{14. Points bloquants}{clôture}
\dire{$\bullet$ Fait~: 4 pistes, chapitre identifiabilité réécrit, mémoire en \n{5 parties}, \n{13 fiches}.}
\dire{$\bullet$ En attente d'autres~: les 2 relectures, le registre, la posture capital (avec Caroline).}
\dire{$\bullet$ \n{\og Ne jamais présenter comme mesuré ce qui est posé, ni comme calibré ce qui n'est que corroboré.\fg}}

\end{multicols}

\vspace{-4pt}
\begin{alerte}
\footnotesize
\textbf{Les trois questions à anticiper.}\\[3pt]
\textbf{\og Ton $z=+3{,}9$, il vaut quoi ?\fg} $\rightarrow$ Ce que vaut un corpus de 7 documents codés par une seule
personne. Le signe est net, le protocole écrit, mais je ne revendique que la structure d'ensemble. C'est pour ça
que je demande le second codeur.\\[3pt]
\textbf{\og Pourquoi tu ne calibres pas $W$ ?\fg} $\rightarrow$ C'est démontré impossible sur ces données, et c'est
un résultat, pas un aveu. Ce que je peux faire, c'est spécifier la donnée qui le permettrait et chiffrer ce qu'elle
vaut. \emph{(Si on insiste~: la vraisemblance est plate en $A$, deux valeurs opposées donnent la même loi. Il faut
d'autres données, pas plus de données.)}\\[3pt]
\textbf{\og Et le SCR, il vaut combien ?\fg} $\rightarrow$ Pas de niveau tant que la posture n'est pas tranchée avec
Caroline, point prédictif ou borne prudente. Les deux sont chiffrés, l'écart est lui-même un résultat. Ce que
j'affirme aujourd'hui, c'est l'ordre des piliers, et il ne bouge sous aucune hypothèse testée.
\end{alerte}

\vspace{2pt}
\begin{alerte}
\footnotesize
\textbf{Trois pièges à éviter.}
\textbf{(i)} Ne pas dire que $W$ \emph{est} une chaîne de Markov~: c'est un noyau de contagion, j'emprunte
l'algèbre de la réversibilité comme lecture, pas comme hypothèse.
\textbf{(ii)} Ne pas dire \og $P=S+A$\fg~: la décomposition porte sur le flux stationnaire pondéré par $\pi$~; mon
$S=(W+W^{\top})/2$ en est le cas uniforme, écart mesuré $0{,}069$.
\textbf{(iii)} Ne pas dire que la martingale \og ne voit rien\fg~: elle voit très bien $S$, elle est aveugle à $A$.\\[3pt]
\textbf{Si Hugo ne trouve personne pour le codage} $\rightarrow$ proposer \textbf{Coraline} ou \textbf{Nathanael},
sujets adjacents, ne connaissent pas l'hypothèse. Mieux vaut un nom qu'un \og je vais voir\fg.\\[3pt]
\textbf{Après le call} $\rightarrow$ envoyer les 2 PDF le jour même, avec une \textbf{date de retour} explicite.
\end{alerte}

\end{document}
